
\section{Appendix dedicated to the optimal algorithms for $\textsc{Uniform}$ and $\textsc{Bernoulli}$ functions}
\label{app:lumbroso}

\noteAnt{mettre un peu de blabla}

\begin{algorithm}[htb]
\caption{Generating a random integer} 
\label{algo:tableXX}
	\begin{algorithmic}[1]
        \renewcommand{\algorithmicrequire}{\textbf{Input:}}
        \renewcommand{\algorithmicensure}{\textbf{Output:}}
        \Require $n$, the upper limit (excluded)
        \Ensure  A random number between $0$ and $n-1$.
        \Function{rand\_int}{$n$}
            \While{True}
                \State $res = 0$
                \State $bits = n.bit\_length()$ \Comment{Return the number of necessary bits to represent the integer in binary}
                \For{$i=0 \text{ to } bits$}
                    \State $res = res * 2 + randint(0, 1)$ \Comment{Add the bit generated}
                \EndFor
                \If{$res < n$}
                    \State \Return $res$
                \EndIf
            \EndWhile
        \EndFunction
	\end{algorithmic}   
 \end{algorithm}

\begin{algorithm}[htb]
\caption{Random sampling in $\bar{T}$} 
\label{algo:tableXY}
	\begin{algorithmic}[1]
        \renewcommand{\algorithmicrequire}{\textbf{Input:}}
        \renewcommand{\algorithmicensure}{\textbf{Output:}}
        \Require $\bar{T}$, $\cs, \bar{q}$  previously defined.
        \Ensure A random object from $\bar{T}$
        \Function{RandomGeneration}{$\bar{T}$, $\cs, \bar{q}$}
            \State $q1 = \bar{q} - 1$
            \State $q2 = \bar{q} - 2$
            \While{True}
                \State $r = \texttt{rand\_int}(\bar{q})$
                \State $box = \bar{T}[r]$
                \If{$0 \leq r \leq q2$} \Comment{First part of the table $\bar{T}$}
                    \State $r' = \texttt{rand\_int}(cs)$
                    \If{$r' \leq box[0]$}
                        \State \Return $box[1]$
                    \Else
                        \State \Return $box[2]$
                    \EndIf
                \ElsIf{$r == q1$} \Comment{Last box of the table $\bar{T}$ : there is potentially a virtual object}
                    \State $r' = \texttt{rand\_int}(cs)$
                    \If{$r' \leq box[0]$}
                        \State \Return $box[1]$
                    \EndIf
                \EndIf
            \EndWhile
        \EndFunction
    \end{algorithmic}
\end{algorithm}




\section{Appendix dedicated to the Combinatorial Chemistry case study}
\label{app:chemistry}
\noteAnt{A nettoyer, des choses ont déjà été dites dans corps du papier}

\paragraph{Introduction and Problem Statement}
Weighted random sampling is a fundamental operation in scientific computing. It consists in selecting an item $x_i$ from a set of $n$ candidates $X = \{x_0, x_1, \dots, x_{n-1}\}$ with probability strictly proportional to its integer weight $w_i$:
\begin{equation}
\mathbb{P}(X = x_i) = \frac{w_i}{\text{total}}
\end{equation}

\noteAnt{il y a qq chose de commenter qui ne compile pas, mais est iportant à mettre en place}
% counter = ncs.ApproximateCounter(show_progress=False)\\
% space = ncs.SearchSpace("C:4 H:1 O:2 N:3 F:1") ## ce sont les valences correctes présentes dans la nature -- vallence == nb de liaisons // dans l'exemple précédent on autorisait des valences + importantes\\
% mols = ncs.random_sample(
%     counter, space, natoms=32, nmols=1, selection=ncs.Q("C + H + O + N + F = 32")
% )}
% $N = 20\,484$; $Z = 467734548828719531605240429605504136894484614517198 \sim 1.25 \cdot 2^{168} \sim 4.7 \cdot 10^{50}$.


\paragraph{Floating-Point Arithmetic Bias: Why \texttt{random.uniform} is Not Uniform}
In Python, the standard implementation for cumulative weighted sampling draws a continuous random variable $r \sim \mathcal{U}[0, \text{total})$ and searches for the corresponding cumulative range. Within Python's standard library (\texttt{random.choices}), sampling relies on calling \texttt{random.uniform(0, total)} or \texttt{random() * total}. The weighted sampling step (5) is implemented in \texttt{nablachem/space/algos.py} as shown in Figure~4:

% [FIGURE 4: Chemistry sampling code snippet (_weighted_choose)]

Although \texttt{total} is accumulated using Python's arbitrary-precision integers (\texttt{int}), evaluating \texttt{random.uniform(0, total)} implicitly converts this integer mass into an IEEE 754 double-precision float (\texttt{float64}). An IEEE 754 \texttt{float64} value allocates 52 explicit mantissa bits plus 1 implicit bit, granting $M = 2^{53}$ distinct representable values in $[0, 1)$.

As long as $\text{total} \le 2^{53}$, every integer index possesses an exact floating-point representation. However, as soon as $\text{total} > 2^{53}$, the step size between consecutive representable floating-point values becomes strictly greater than $1$:
\begin{itemize}
    \item Between $2^{53}$ and $2^{54}$, the step size is $2^1 = 2$: only even integers can be reached.
    \item Between $2^{54}$ and $2^{55}$, the step size is $2^2 = 4$: only multiples of 4 are reachable.
    \item For $\text{total} \sim 2^{100}$, the step size between consecutive draws leaps to $2^{47} \approx 1.4 \times 10^{14}$.
\end{itemize}

Let $M = 2^{\text{bits}}$ be the number of representable values for a uniform floating-point number $f \in [0, 1)$, and $\text{total}$ be the total weight of the distribution:
\begin{itemize}
    \item \textbf{Regime $M \gg \text{total}$ (sufficient precision):} Discretization error is limited to minor rounding noise.
    \item \textbf{Regime $M \ll \text{total}$ (large-scale real-world case, e.g., $M = 2^{53}, \text{total} \sim 2^{100}$):} The resolution step $\text{total}/M$ exceeds the individual weights of many candidate items. An item whose weight is smaller than this grid resolution incurs a selection probability that becomes \emph{structurally zero}, regardless of its non-zero theoretical probability.
\end{itemize}

To quantify this deviation exactly without resorting to expensive Monte Carlo simulations, an $O(n)$ exact arithmetic algorithm calculates the effective probability distribution reachable for a target bit precision. By evaluating how many grid points $f = k/M$ (projected as $f \times \text{total}$) land within each item's cumulative interval, the algorithm computes effective sampling weights using integer division with ceiling, completely independent of the magnitude of $\text{total}$.

% [FIGURE 5: Stretching and shrinking of a finite-precision random source (M=16)]

Figure~5 illustrates these discretization mechanics using 4-bit precision ($M = 16$, center axis):
\begin{itemize}
    \item \textbf{Stretched regime ($M < \text{total}$, top axis, blue):} The 16 available values are stretched over 22 target integer slots. Because the resolution step ($22/16 = 1.375$) is strictly greater than 1, covering every target integer is impossible. Gaps remain: items falling entirely into these gaps have a zero probability of being chosen.
    \item \textbf{Shrunk regime ($M > \text{total}$, bottom axis, orange):} The 16 values are compressed over 10 target integer slots. The average step ($16/10 = 1.6$) allows every target to be reached. However, since 10 does not divide 16, grid points stack unevenly, assigning two points to certain targets and one to others, inducing a modulo bias.
\end{itemize}

\paragraph{Advantages of the Integer ALIAS Method in Arbitrary Precision}
The ALIAS method using arbitrary-precision integers provides an optimal solution to address both precision errors during sampling and execution time bottlenecks.

First, this method guarantees complete mathematical exactness. By operating exclusively with arbitrary-precision integers, it preserves the exact weight of every item regardless of scale. The resulting draws remain strictly unbiased and faithful to the theoretical distribution.

Second, it provides significant speedups. Linear scanning methods iterate through weights sequentially, incurring an $O(n)$ sampling time per draw. Across millions of trials, this becomes prohibitive. Our ALIAS approach resolves this by building a lookup table once during a preprocessing step, after which every subsequent draw executes in $O(1)$ constant time.

In summary, arbitrary-precision integer ALIAS eliminates floating-point conversion errors while enabling instant random sampling, making it ideal for massive combinatorial spaces.




\section{Lumbroso's sampling algorithms}


\noteAnt{Algo de Lumbroso \cite{Lumbroso13}}
\begin{algorithm}[H]
\caption{Generating a random integer or a Bernoulli variable} 
\label{algo:RandInt_v_AI}
    \begin{tabular}{l l}
        \begin{minipage}{0.5\textwidth}
	        \begin{algorithmic}[1]
                \renewcommand{\algorithmicrequire}{\textbf{Input:}}
                \renewcommand{\algorithmicensure}{\textbf{Output:}}
                \Require $n$, the upper limit (excluded)
                \Ensure  A uniform random integer between $0$ and $n-1$.
                \Function{\textsc{Uniform}}{$n$}
                    \State $k = \lceil \log_2(n) \rceil$
                    \State $bound = 2^k$
                    \State $x = \texttt{Rand\_Bits}(k)$
                    \While{$bound < n$ \textbf{or} $x \geq n$}
                        \If{$bound \geq n$}
                            \State $bound = bound - n$
                            \State $x = x - n$
                        \EndIf
                        \State $bound = 2 \cdot bound$
                        \State $x = 2 \cdot x + \texttt{Rand\_Bits}(1)$
                    \EndWhile
                    \State \Return $x$
                \EndFunction
	        \end{algorithmic}
        \end{minipage}
        &
        \begin{minipage}{0.5\textwidth}
            \begin{algorithmic}[1]
                \renewcommand{\algorithmicrequire}{\textbf{Input:}}
                \renewcommand{\algorithmicensure}{\textbf{Output:}}
                \Require $p, q$ with $0 \leq p \leq q$
                \Ensure  A Bernoulli variable equal to $1$ with probability $p/q$.
                \Function{\textsc{Bernoulli}}{$p, q$}
                    \If{$p = q$}
                        \State \Return $1$
                    \EndIf
                    \State $p = 2 \cdot p$
                    \While{$p \neq q$}
                        \If{$p > q$}
                            \State $y = 1$
                            \State $p = p - q$
                        \Else
                            \State $y = 0$
                        \EndIf
                        \If{\texttt{Rand\_Bits}(1) == 1}
                            \State \Return $y$
                        \EndIf
                        \State $p = 2 \cdot p$
                    \EndWhile
                    \State \Return \texttt{Rand\_Bits}(1)
                \EndFunction
            \end{algorithmic}
        \end{minipage}
    \end{tabular}
    \begin{itemize}
        \footnotesize 
        \item $\texttt{Rand\_Bits}(k)$: it returns an integer whose base-2 representation is built with a sequence of $k$ unbiased random bits
    \end{itemize}
 \end{algorithm}
